% !TEX root = ../my-thesis.tex
%


\lstset{breaklines=true, breakatwhitespace=false}

\chapter{Praktische Modellierung ausgewählter Workflow-Patterns}
\label{sec:practical}

Dieses Kapitel dokumentiert die praktische Umsetzung zweier Workflow-Patterns in den fünf
betrachteten Systemen. In Abschnitt~\ref{sec:practical:setup} werden Zielsetzung,
Benchmark-Skripte, Testumgebung und die Einrichtung der Systeme beschrieben. Die Abschnitte
\ref{sec:practical:pipeline} und \ref{sec:practical:scattergather} behandeln jeweils ein
Pattern und gehen dabei für jedes System einzeln durch, wie die Umsetzung aussieht. Am Ende
beider Abschnitte stehen vergleichende Beobachtungen. Abschnitt~\ref{sec:practical:overall}
fasst die systemübergreifenden Erkenntnisse zusammen.

\section{Zielsetzung und Testumgebung}
\label{sec:practical:setup}

\subsection{Zielsetzung und Vorgehensweise}
\label{sec:practical:setup:goal}

Ziel der praktischen Modellierung ist es, zwei ausgewählte Workflow-Patterns, das
\emph{Pipeline-Pattern} und das \emph{Scatter-Gather-Pattern} (bestehend aus den beiden Mustern
Parallel-Split und Synchronization), in allen fünf betrachteten Systemen umzusetzen und dabei die
systemspezifischen Unterschiede in der Workflow-Komposition, der Datenweitergabe und der
Ausführungssteuerung sichtbar zu machen. Beide Patterns wurden in Kapitel~\ref{sec:basics}
eingeführt. Für beide Patterns wird dieselbe Grundrechenlogik verwendet, sodass die Systeme unter
vergleichbaren Bedingungen betrachtet werden können und sich die beobachteten Unterschiede auf die
Workflow-Modellierung und nicht auf die Berechnung zurückführen lassen.

Um die Systeme nicht nur einzeln zu beschreiben, sondern systematisch zu vergleichen, wird die
Umsetzung in jedem System entlang derselben Kriterien betrachtet, die sich an dem in dieser Arbeit
entwickelten Bewertungs- und Vergleichsrahmen orientieren. Diese Kriterien werden im Folgenden
nicht stur Punkt für Punkt abgearbeitet, sondern fließen in die Beschreibung der einzelnen Systeme
ein:

\begin{itemize}
  \item \emph{Ausdruck des Patterns}: Wie wird die Struktur des Patterns, also die sequenzielle
    Abfolge bzw. die Verzweigung und Zusammenführung, in der jeweiligen Beschreibungssprache
    modelliert?
  \item \emph{Modellierungsaufwand}: Wie viel Beschreibungs- und Konfigurationsaufwand erfordert
    die Umsetzung?
  \item \emph{Sichtbarkeit von Datenfluss und Abhängigkeiten}: Wie entstehen die Abhängigkeiten
    zwischen den Schritten, und wie sichtbar ist der dateibasierte Datenfluss?
  \item \emph{Ausführung und Steuerung}: Wie wird die Ausführung gesteuert, und wie lässt sich das
    beobachtete Laufzeitverhalten, etwa die Parallelität, nachweisen?
  \item \emph{Trennung von Workflow-Logik und Ausführungsumgebung}: Wie stark ist die
    Workflow-Beschreibung von der konkreten Ausführungsumgebung entkoppelt?
  \item \emph{Verständlichkeit aus Einsteigerperspektive}: Wie zugänglich ist die Umsetzung für
    Personen ohne Vorerfahrung mit dem jeweiligen System?
\end{itemize}

Die folgenden Umsetzungsabschnitte orientieren sich an diesem Beobachtungsrahmen. Für jedes System
wird daher betrachtet, wie das Pattern beschrieben wird, wie Datenfluss und Abhängigkeiten
entstehen, wie die Ausführung erfolgt und welche Informationen zur Nachvollziehbarkeit zur
Verfügung stehen. Dadurch sollen die Systeme nicht nur beschrieben, sondern entlang derselben
Punkte vergleichbar beobachtet werden.

\subsection{Benchmark-Skripte}
\label{sec:practical:setup:scripts}

Beiden Patterns liegt dieselbe, bewusst einfach gehaltene Rechenlogik zugrunde, damit der Fokus
auf der Workflow-Modellierung und nicht auf der Berechnung selbst liegt. Den Kern bilden vier
Python-Skripte. \texttt{generate\_input.py} erzeugt zufällige Ganzzahlen im Bereich von 1 bis 30
und schreibt sie zeilenweise in eine Datei (beim Pipeline-Pattern zehn Werte). \texttt{preprocess.py}
liest diese Datei wieder ein und filtert alle Werte kleiner als 10 heraus. \texttt{compute.py}
berechnet Anzahl, Summe und Mittelwert der verbliebenen Werte, und \texttt{postprocess.py} fasst
das Ergebnis in einer lesbaren \texttt{summary.txt} zusammen.

Für das Pipeline-Pattern wurden diese Skripte in Nextflow, Pegasus und Merlin direkt aufgerufen.
Für StreamFlow mussten sie angepasst werden, da Eingabe- und Ausgabepfade in der
containerisierten Ausführungsumgebung als Kommandozeilenargumente übergeben werden. In AiiDA
wurde dieselbe Rechenlogik innerhalb der \texttt{WorkChain}-Methoden umgesetzt, ohne externe
Skriptaufrufe.

Für das Scatter-Gather-Pattern wurde diese Grundlogik um zwei Skripte erweitert. \texttt{split.py}
teilt die vorbereiteten Werte auf vier etwa gleich große Teildateien
\texttt{chunk\_1.txt} bis \texttt{chunk\_4.txt} auf (Scatter), und \texttt{aggregate.py} liest die
vier Teilergebnisse ein, summiert Anzahl und Summe über alle Chunks und berechnet daraus den
Gesamtmittelwert (Gather). Zusätzlich wurde die Eingabemenge auf 100 Werte erhöht, damit die
Aufteilung auf mehrere Chunks überhaupt sinnvoll ist, und sämtliche Skripte nehmen ihre Ein- und
Ausgabepfade über \texttt{sys.argv} entgegen, da die vier parallelen \texttt{compute}-Aufrufe
jeweils unterschiedliche Dateien verarbeiten. In dieser parametrisierten Form wurden die Skripte
in allen Systemen unverändert verwendet, außer in AiiDA. In AiiDA wurde dieselbe Rechenlogik innerhalb der WorkChain-Implementierung nachgebildet. Die vollständigen Skripte sind in Anhang~\ref{app:scripts}
(Pipeline) und Anhang~\ref{app:sg:scripts} (Scatter-Gather) abgebildet.

\subsection{Testumgebung}
\label{sec:practical:setup:environment}

Alle Implementierungen wurden lokal auf einem MacBook Pro mit Apple-Silicon-Prozessor (arm64)
unter macOS ausgeführt. Als Paketverwaltung kamen je nach System Homebrew, conda und pip zum
Einsatz. Es handelt sich damit nicht um Läufe auf einem echten HPC-Cluster, sondern um eine rein
lokale Ausführung auf einem einzelnen Rechner. Die Nähe zum HPC-Umfeld ergibt sich nicht aus der
Rechenlast, sondern aus der Struktur der Workflows, bei der Daten vorbereitet, aufgeteilt,
unabhängig verarbeitet und wieder zusammengeführt werden. Auf diese Einschränkung wird in
Abschnitt~\ref{sec:practical:overall} noch einmal eingegangen.
Tabelle~\ref{tab:practical:setup} gibt einen Überblick über den Einrichtungsaufwand je System.

\subsection{Einrichtung der Systeme}
\label{sec:practical:setup:installation}

Schon vor der eigentlichen Modellierung zeigte sich ein deutlicher Unterschied zwischen den
Systemen, nämlich beim Aufwand für Installation und Einrichtung. Dieser Punkt ist für die Praxis
nicht nebensächlich, da ein hoher Einrichtungsaufwand für jemanden, der nur einen kleinen Workflow
ausprobieren möchte, bereits eine Einstiegshürde darstellt, noch bevor der erste Workflow läuft.
Die folgenden Beschreibungen beziehen sich auf Probleme, die nicht an einem Pattern, sondern an
der Installation und Konfiguration der jeweiligen Infrastruktur lagen. Die genauen Befehle sind
nicht abgedruckt, da sie für das Verständnis der Patterns nicht zentral sind.

Nextflow~25.10.4 wurde über Homebrew installiert und benötigt keine externen Dienste. Bei der
Einrichtung traten keine Probleme auf, sodass unmittelbar mit der Modellierung begonnen werden
konnte. Pegasus~5.1.2 erforderte ein hohes Vorbereitungsaufwand. Neben Pegasus selbst musste
HTCondor~23.1.0 installiert werden, das die Ausführung übernimmt, und die Pegasus-Python-Bibliothek
über die Umgebungsvariable \texttt{PYTHONPATH} bekannt gemacht werden. Beim ersten Versuch schlug
die Planung fehl, weil \texttt{CONDOR\_CONFIG} nicht gesetzt war und HTCondor seine Konfiguration
nicht finden konnte. Die Lösung bestand darin, vor jedem Aufruf ein Skript zu sourcen, das diese
Variable setzt.

Merlin~1.13.0 ließ sich an sich problemlos installieren, startete danach aber zunächst nicht. Der
Grund war, dass die aktuelle Version von \texttt{setuptools} das intern benötigte Modul
\texttt{pkg\_resources} nicht mehr mitliefert, das Merlin voraussetzt. Nach einem Downgrade auf
eine Version kleiner als 81 lief alles. Als Hintergrunddienste verwendet Merlin Redis als
Nachrichtenspeicher und Celery für die Verteilung der Aufgaben, was sich über einen
Docker-Container bereitstellen ließ. StreamFlow~0.1.6 benötigt Python in Version 3.12, da ein
intern genutztes Modul in neueren Python-Versionen entfernt wurde, sowie Docker für die Ausführung
in Containern. Beides war für sich genommen kein großes Problem, allerdings trat später bei der
Ausführung ein Kompatibilitätsproblem mit Docker auf, das im Abschnitt zum Pipeline-Pattern
genauer beschrieben wird.

Den mit Abstand höchsten Aufwand hatte AiiDA~2.8.0. Hier reicht es nicht, das Werkzeug selbst zu
installieren, sondern es setzt eine PostgreSQL-Datenbank voraus, in der es seine Daten ablegt,
sowie RabbitMQ als Nachrichtendienst, über den die Aufgaben an einen Hintergrundprozess vermittelt
werden. Zusätzlich mussten eine Datenbank angelegt, ein Profil eingerichtet und ein Daemon
gestartet werden, der die Workflows im Hintergrund ausführt. Bei jedem Lauf erschien zudem eine
Warnung, dass die installierte RabbitMQ-Version offiziell nicht unterstützt wird. Für die kurzen
Testworkflows hatte das keine Auswirkung, für einen produktiven Einsatz wäre aber eine
unterstützte Version erforderlich. Dieser hohe Aufwand zahlt sich an anderer Stelle aus, weil
AiiDA dadurch die gesamte Daten- und Prozesshistorie speichert, worauf später noch eingegangen
wird.

\begin{table}[ht]
  \caption{Benötigte externe Dienste und Einrichtungsaufwand je System}
  \label{tab:practical:setup}
  \centering
  \begin{tabular}{lll}
    \hline
    \textbf{System} & \textbf{Externe Dienste} & \textbf{Aufwand} \\
    \hline
    Nextflow   & keine                              & sehr gering \\
    Merlin     & Redis, Celery                      & gering \\
    StreamFlow & Docker                             & mittel \\
    Pegasus    & HTCondor                           & hoch \\
    AiiDA      & PostgreSQL, RabbitMQ, Daemon       & sehr hoch \\
    \hline
  \end{tabular}
\end{table}

\section{Pipeline-Pattern}
\label{sec:practical:pipeline}

\subsection{Workflow-Beschreibung}
\label{sec:practical:pipeline:description}

Das Pipeline-Pattern (vgl. Abschnitt~\ref{sec:basics}) wurde als strikt sequenzieller
Vier-Schritt-Workflow umgesetzt, bei dem jeder Schritt erst beginnt, wenn der vorherige
abgeschlossen ist, und dessen Ausgabe als Eingabe des nächsten dient:

\[
  \texttt{generate\_input} \;\rightarrow\; \texttt{preprocess} \;\rightarrow\;
  \texttt{compute} \;\rightarrow\; \texttt{postprocess}
\]

Die vier Benchmark-Skripte aus Abschnitt~\ref{sec:practical:setup:scripts} bilden die
Rechengrundlage. Das Pattern eignet sich als erster Vergleichsworkflow, weil es sich in allen
fünf Systemen umsetzen lässt und bereits grundlegende Konzepte wie Abhängigkeiten, Datenweitergabe
und die sequenzielle Abfolge sichtbar macht. Die folgenden Abschnitte beschreiben die Umsetzung in
den einzelnen Systemen. Auch wenn alle Systeme denselben Ablauf abbilden, unterscheiden sie sich
deutlich darin, wie der Workflow beschrieben wird, wie die Reihenfolge entsteht und wie die Daten
weitergegeben werden.

\subsection{Umsetzung in Nextflow}
\label{sec:practical:pipeline:nextflow}

\paragraph{Workflow-Komposition.}
Der Workflow wird in der Datei \texttt{main.nf} in der Nextflow-eigenen Sprache DSL2 beschrieben.
Jeder Rechenschritt ist ein \texttt{process}-Block mit seinen Ein- und Ausgaben und ein
\texttt{workflow}-Block setzt die vier Prozesse zusammen. Dieser \texttt{workflow}-Block bildet den
Kern der Beschreibung und besteht im Wesentlichen aus den vier Prozessaufrufen:

\begin{lstlisting}[caption={workflow-Block aus main.nf (Nextflow, Pipeline)},label={lst:pipe:nf:workflow}]
workflow {
    raw_ch      = GENERATE_INPUT()
    prepared_ch = PREPROCESS(raw_ch)
    result_ch   = COMPUTE(prepared_ch)
    summary_ch  = POSTPROCESS(result_ch)
}
\end{lstlisting}

Die vollständige Datei einschließlich der \texttt{process}-Blöcke ist in
Anhang~\ref{app:nextflow:mainnf} abgebildet.

\paragraph{Datenfluss und Abhängigkeiten.}
Die Reihenfolge der vier Schritte ergibt sich allein daraus, dass die Ausgabe eines Prozesses über
einen Channel als Eingabe des nächsten verwendet wird. Ein eigenes Schlüsselwort wie
\texttt{depends} ist nicht nötig. So wird etwa \texttt{prepared\_ch}, die Ausgabe von
\texttt{PREPROCESS}, direkt an \texttt{COMPUTE} weitergereicht. Eine als \texttt{path} deklarierte
Datei stellt Nextflow automatisch per Symlink im Arbeitsverzeichnis des jeweiligen Prozesses
bereit, sodass die Skripte sie direkt öffnen können. Da jeder Prozess in einem eigenen
Unterverzeichnis unterhalb von \texttt{work/} arbeitet, mussten die Skripte über einen absoluten
Pfad aufgerufen werden, wofür Nextflow die Variable \texttt{projectDir} bereitstellt.

\paragraph{Ausführung und Umgebung.}
Der Workflow lief im lokalen Executor fehlerfrei durch, wobei alle vier Schritte jeweils als
abgeschlossen gemeldet wurden. Externe Dienste oder Hintergrundprozesse sind dafür nicht
erforderlich, was Nextflow in der Handhabung besonders unkompliziert macht. Der selbe Workflow ließe sich grundsätzlich auch an einen HPC-Scheduler
übergeben, ohne die Beschreibung in \texttt{main.nf} zu ändern.

\paragraph{Nachvollziehbarkeit und Einordnung.}
Nach dem Lauf lässt sich anhand der isolierten Arbeitsverzeichnisse nachvollziehen, welcher Schritt
welche Datei erzeugt hat, was die Schritte klar voneinander trennt und reproduzierbar macht. Aus
Einsteigersicht ist die Beschreibung sehr gut lesbar, lediglich der Umgang mit \texttt{projectDir}
und den isolierten Verzeichnissen erfordert etwas Einarbeitung. Im Vergleich zu den anderen
Systemen ist die Umsetzung in Nextflow die kompakteste, da der gesamte Ablauf in einer einzigen,
kurzen Datei steht.

\subsection{Umsetzung in Pegasus}
\label{sec:practical:pipeline:pegasus}

\paragraph{Workflow-Komposition.}
Bei Pegasus wird der Workflow nicht in einer eigenen Sprache, sondern über die Pegasus Python API
programmatisch aufgebaut. Das Skript \texttt{build\_workflow.py} registriert die Skripte zunächst
als sogenannte Transformations und legt anschließend die Jobs an, aus denen die eigentlichen
Workflow-Dateien erzeugt werden. Die vollständige Datei ist in Anhang~\ref{app:pegasus:buildwf}
abgebildet.

\paragraph{Datenfluss und Abhängigkeiten.}
Anders als bei Nextflow wird die Reihenfolge nicht über die Aufrufreihenfolge, sondern über die
Ein- und Ausgabedateien der Jobs festgelegt. Jeder Job deklariert mit \texttt{add\_inputs} und
\texttt{add\_outputs}, welche Dateien er liest und schreibt:

\begin{lstlisting}[language=Python,caption={Ein- und Ausgaben eines Jobs (Pegasus, Pipeline)},label={lst:pipe:pg:io}]
compute_job = Job("compute")
compute_job.add_inputs(prepared_file)
compute_job.add_outputs(result_file, stage_out=False, register_replica=False)
\end{lstlisting}

Erzeugt ein Job eine Datei, die ein anderer als Eingabe nutzt, leitet Pegasus daraus automatisch
eine Kante im gerichteten azyklischen Graphen (DAG) ab. Der Datenfluss ist dadurch sehr explizit,
verlangt aber, dass für jeden Job die beteiligten Dateien einzeln angegeben werden.

\paragraph{Ausführung und Umgebung.}
Ein kennzeichnendes Merkmal von Pegasus ist die Trennung von Planung und Ausführung. Der Befehl
\texttt{pegasus-plan} übersetzt den abstrakten Workflow in einen konkreten HTCondor-DAG und löst
dabei Abhängigkeiten und Dateizugriffe auf, erst \texttt{pegasus-run} übergibt ihn an HTCondor,
das die Jobs ausführt. Als Datenkonfiguration wurde \texttt{sharedfs} verwendet, da im lokalen
Setup Submit-Seite und HTCondor-Ausführungsort auf dieselben Verzeichnisse zugreifen. Dadurch ist
kein Deployment eines Pegasus-Worker-Pakets nötig, und die Schritte tauschen ihre Dateien über das
gemeinsame Dateisystem aus, während Pegasus weiterhin die Abhängigkeitssteuerung, die Ausführung
über HTCondor sowie Stage-out, Registrierung und Cleanup übernimmt. Der Lauf war erfolgreich. Die
Workflow-Wall-Time lag bei rund zwei Minuten und siebzehn Sekunden, die kumulierte Job-Laufzeit
bei etwa 28 Sekunden. Diese Gesamtzeit ist nicht als reine Berechnungs- oder Ausführungszeit zu
verstehen, da sie neben den sehr kurzen Python-Schritten auch Planung, HTCondor-Queueing,
DAGMan-Steuerung sowie Stage-out, Registrierung und Cleanup einschließt.

\paragraph{Nachvollziehbarkeit und Einordnung.}
Nach dem Lauf lässt sich der erzeugte DAG einsehen, der die Abhängigkeiten zwischen den Schritten
deutlich macht. Bei so kurzen Tasks und einem lokalen Ein-Slot-HTCondor-Aufbau prägen Verwaltungs-
und Schedulinganteile den gesamten Makespan, weshalb sich aus der gemessenen Zeit keine Aussage
über das Verhalten großer HPC-Produktivsysteme ableiten lässt, für die Pegasus gedacht ist. Die
Trennung von abstraktem Workflow und konkreter Ausführung ist klar gegeben, da derselbe Workflow
grundsätzlich auch für andere Sites geplant werden könnte. Aus Einsteigersicht ist der Aufwand
durch die Kataloge und die Trennung von Planung und Ausführung deutlich höher als bei Nextflow.

\subsection{Umsetzung in Merlin}
\label{sec:practical:pipeline:merlin}

\paragraph{Workflow-Komposition.}
Merlin beschreibt den Workflow in einer YAML-Datei \texttt{pipeline.yaml} im \texttt{study}-Block,
in dem die vier Schritte nacheinander aufgeführt sind. Die vollständige Datei ist in
Anhang~\ref{app:merlin:yaml} abgebildet.

\paragraph{Datenfluss und Abhängigkeiten.}
Die Reihenfolge wird über das Schlüsselwort \texttt{depends} angegeben, was den Ablauf gut lesbar
macht. Anders als bei Nextflow werden die Dateien jedoch nicht automatisch weitergegeben, sondern
zu Beginn jedes Schritts von Hand aus dem Arbeitsverzeichnis des vorherigen kopiert:

\begin{lstlisting}[caption={depends und manuelles Kopieren (Merlin, Pipeline)},label={lst:pipe:me:depends}]
  - name: compute
    run:
      depends: [preprocess]
      cmd: |
        cp $(preprocess.workspace)/prepared_input.txt .
        python3 $(SPECROOT)/scripts/compute.py
\end{lstlisting}

Die Abhängigkeit ist damit über \texttt{depends} unmittelbar ablesbar, der Datenfluss ist aber
weniger automatisch als bei den anderen Systemen, weil die Dateien explizit über
\texttt{\$(schritt.workspace)} kopiert werden müssen.

\paragraph{Ausführung und Umgebung.}
Merlin kennt zwei Ausführungsarten. Im lokalen Modus werden die Schritte direkt im aktuellen
Prozess abgearbeitet, im Worker-Modus legt Merlin die Schritte als Aufgaben in Redis ab, die von
Celery-Workern bearbeitet werden. Beim Pipeline-Pattern brachte der Worker-Modus noch keinen
Vorteil, weil die Schritte ohnehin nacheinander laufen; dieser zeigt sich erst beim
Scatter-Gather-Pattern in Abschnitt~\ref{sec:practical:scattergather}.

\paragraph{Nachvollziehbarkeit und Einordnung.}
Da jeder Schritt sein eigenes Workspace-Verzeichnis besitzt, lässt sich nach dem Lauf gut
nachvollziehen, welche Dateien wo entstanden sind. Die Verteilung über Redis und Celery gehört zur
Ausführungsumgebung und ist von der Workflow-Beschreibung getrennt. Aus Einsteigersicht ist die
YAML-Beschreibung gut lesbar und die \texttt{depends}-Angaben sind selbsterklärend, das manuelle
Kopieren der Dateien ist jedoch ein zusätzliches Konzept, das man kennen muss. Im Vergleich liegt
Merlin damit zwischen der sehr automatischen Lösung von Nextflow und dem expliziten Ansatz von
Pegasus.

\subsection{Umsetzung in StreamFlow}
\label{sec:practical:pipeline:streamflow}

\paragraph{Workflow-Komposition.}
StreamFlow trennt die Workflow-Logik und die Ausführungsumgebung in zwei Dateien. Die fachliche
Beschreibung steht in \texttt{workflow.cwl} in der \emph{Common Workflow Language} (CWL), die
Ausführungsumgebung in \texttt{streamflow.yml}, und eine dritte Datei \texttt{config.yml} übergibt
die Skripte. Da die Schritte in Docker-Containern laufen, nehmen die Skripte ihre Pfade als
Argumente entgegen, die in der CWL über \texttt{inputBinding} übergeben werden. Die vollständigen
Dateien sind in Anhang~\ref{app:streamflow:cwl} abgebildet.

\paragraph{Datenfluss und Abhängigkeiten.}
Die Abhängigkeiten entstehen implizit, indem die Ausgabe eines Schritts als Eingabe des nächsten
referenziert wird. So bezieht der Schritt \texttt{preprocess} seine Eingabe direkt aus der Ausgabe
von \texttt{generate\_input}:

\begin{lstlisting}[caption={Implizite Abhängigkeit über eine CWL-Referenz (StreamFlow, Pipeline)},label={lst:pipe:sf:ref}]
    in:
      script:   preprocess_script
      raw_file: generate_input/raw_file
    out: [prepared_file]
\end{lstlisting}

Der Datenfluss ist über solche Referenzen vollständig nachvollziehbar. StreamFlow kopiert die
Dateien in temporäre Container-Verzeichnisse und das Endergebnis zurück auf das lokale Dateisystem.

\paragraph{Ausführung und Umgebung.}
Die Schritte werden in Docker-Containern ausgeführt. Dabei trat das auffälligste Problem der
gesamten Testreihe auf: StreamFlow setzte intern ein Docker-Flag, das neuere Docker-Versionen nicht
mehr unterstützen, woran das Auslesen einer internen Ausgabe scheiterte. Erst nach einer Anpassung
in der internen Quelldatei \texttt{container.py} lief der Workflow durch. StreamFlow war damit das
einzige der fünf Systeme, bei dem ein Eingriff in den Quellcode des Werkzeugs selbst nötig war.

\paragraph{Nachvollziehbarkeit und Einordnung.}
Das Endergebnis wird als CWL-Output zurückgegeben, und der Datenfluss bleibt über die CWL-Referenzen
nachvollziehbar. Die Trennung von Workflow-Logik und Ausführungsumgebung ist das deutlichste
Merkmal von StreamFlow, da dieselbe CWL ohne Änderung auch für andere Ausführungsmodelle verwendet werden könnte. Aus Einsteigersicht ist die standardisierte CWL gut
nachvollziehbar, der nötige Eingriff in fremden Quellcode stellt jedoch eine Hürde dar,
die über die eigentliche Workflow-Modellierung hinausgeht.

\subsection{Umsetzung in AiiDA}
\label{sec:practical:pipeline:aiida}

\paragraph{Workflow-Komposition.}
In AiiDA wird der Workflow nicht über eine eigene Datei, sondern direkt als
Python-\texttt{WorkChain} umgesetzt. Die einzelnen Schritte sind Methoden der Klasse, und ihre
Reihenfolge wird in \texttt{spec.outline()} festgelegt:

\begin{lstlisting}[language=Python,caption={Festlegung der Schrittreihenfolge (AiiDA, Pipeline)},label={lst:pipe:ai:outline}]
spec.outline(
    cls.generate_input,
    cls.preprocess,
    cls.compute,
    cls.postprocess,
)
\end{lstlisting}

Die vollständige Umsetzung ist in Anhang~\ref{app:aiida:workchain} abgebildet.

\paragraph{Datenfluss und Abhängigkeiten.}
Die Reihenfolge ist in \texttt{spec.outline()} direkt ablesbar. Zwischenergebnisse werden über
einen internen Kontext \texttt{self.ctx} von einer Methode zur nächsten weitergegeben, und die
finalen Ausgaben werden als typisierte AiiDA-Datenknoten wie \texttt{List}, \texttt{Dict} oder
\texttt{SinglefileData} deklariert. Der Datenfluss ist dadurch weniger an einzelne Dateien gebunden
als bei den anderen Systemen und stärker an das Datenmodell von AiiDA.

\paragraph{Ausführung und Umgebung.}
Der Workflow wurde sowohl direkt im aktuellen Prozess als auch über den Daemon ausgeführt, an den
die Aufgabe über RabbitMQ vermittelt wird. Beide Wege liefen erfolgreich durch. Beim Daemon-Lauf
muss der Workflow-Code dem separaten Daemon-Prozess über \texttt{PYTHONPATH} bekannt gemacht werden,
damit dieser die Klasse findet.

\paragraph{Nachvollziehbarkeit und Einordnung.}
Hier zeigt sich die größte Stärke von AiiDA: Es ist das einzige der fünf Systeme, das automatisch
die gesamte Daten- und Prozesshistorie speichert und jeden Schritt sowie die zugehörigen Daten
dauerhaft als Knoten in einer Datenbank ablegt, sodass sich der Ablauf im Nachhinein vollständig
rekonstruieren lässt. Aus Einsteigersicht ist der Code für Personen mit Python-Kenntnissen
verständlich, das Datenmodell und die Einrichtung des Daemons erfordern jedoch zusätzliche
Einarbeitung. Im Vergleich verbindet AiiDA damit eine vertraute Beschreibungsform mit der besten
Nachvollziehbarkeit, allerdings beim höchsten Einrichtungsaufwand.

\subsection{Vergleichende Beobachtungen zum Pipeline-Pattern}
\label{sec:practical:pipeline:comparison}

Alle fünf Systeme konnten das Pipeline-Pattern erfolgreich umsetzen und ausführen.
Tabelle~\ref{tab:pipeline:composition} fasst die Beobachtungen zu den fünf Systemen entlang des
Beobachtungsrahmens zusammen, also Beschreibung, Abhängigkeiten, Datenweitergabe, Ausführung und
Nachvollziehbarkeit.

\begin{table}[ht]
  \caption{Das Pipeline-Pattern je System im Überblick}
  \label{tab:pipeline:composition}
  \centering
  \footnotesize
  \setlength{\tabcolsep}{3pt}
  \renewcommand{\arraystretch}{1.3}
  \begin{tabular}{l p{1.5cm} p{2.0cm} p{2.0cm} p{2.5cm} p{2.2cm}}
    \hline
    \textbf{System} & \textbf{Beschrei\-bung} & \textbf{Abhängig\-keiten} & \textbf{Daten\-weitergabe} & \textbf{Ausführung / Umgebung} & \textbf{Nachvoll\-ziehbarkeit} \\
    \hline
    Nextflow
      & DSL2
      & implizit über Channels
      & auto\-matisch per Symlink
      & lokaler Executor
      & Verzeich\-nisse, Report \\
    Pegasus
      & Python API
      & explizit über Jobs
      & über SharedFS
      & Planung + HTCondor
      & Condor-Logs, DAG \\
    StreamFlow
      & CWL
      & implizit über Referenzen
      & über Container
      & Docker-Container
      & CWL-Outputs, Report \\
    Merlin
      & YAML
      & explizit über \texttt{depends}
      & manuell kopiert
      & lokal oder Worker (Redis, Celery)
      & Work\-spaces je Schritt \\
    AiiDA
      & Python-Klasse
      & explizit über \texttt{outline}
      & über \texttt{ctx} und Knoten
      & direkt oder Daemon (RabbitMQ)
      & speichert Historie \\
    \hline
  \end{tabular}
\end{table}

Schon beim rein sequenziellen Pipeline-Pattern zeigen sich deutliche Unterschiede. Beim Ausdruck
des Patterns und beim Modellierungsaufwand reicht das Spektrum von einer sehr kompakten,
impliziten Modellierung über Channels (Nextflow) und Output-Referenzen (StreamFlow) bis zu
explizit angegebenen Abhängigkeiten über die Python API (Pegasus, AiiDA) oder über \texttt{depends}
(Merlin). Nextflow kam dabei mit dem geringsten, Aiida mit dem höchsten Einrichtungsaufwand aus.

Bei der Ausführung und Steuerung ist bei Pegasus zu beachten, dass die gemessene
Workflow-Wall-Time bei den kurzen Tasks und dem lokalen Ein-Slot-HTCondor-Aufbau stark von
Verwaltungs- und Schedulinganteilen wie Planung, Queueing, Stage-out und Cleanup geprägt ist und
sich daher nicht auf große HPC-Produktivsysteme verallgemeinern lässt. Hinsichtlich der Trennung
von Workflow-Logik und Ausführungsumgebung sind StreamFlow, mit der Trennung von CWL und
\texttt{streamflow.yml}, und Pegasus, mit der Trennung von Planung und Ausführung, am deutlichsten.
StreamFlow war zugleich das einzige System, das einen Eingriff in den eigenen Quellcode erforderte.
Beim Umgang mit den Daten hebt sich AiiDA ab, da es als einziges System die Daten- und
Prozesshistorie automatisch speichert und damit eine besonders gute Nachvollziehbarkeit bietet.
Aus Einsteigersicht ist Nextflow am zugänglichsten, während Pegasus, StreamFlow und AiiDA eine
höhere Einstiegshürde aufweisen.

\section{Scatter-Gather-Pattern}
\label{sec:practical:scattergather}

\subsection{Workflow-Beschreibung}
\label{sec:practical:scattergather:description}

Das Scatter-Gather-Pattern (vgl. Abschnitt~\ref{sec:basics}) verbindet die beiden Muster
Parallel-Split und Synchronization: Ein Datensatz wird in mehrere Teile aufgeteilt, diese Teile
werden unabhängig und parallel verarbeitet, und die Teilergebnisse werden anschließend wieder
zusammengeführt. Für die Umsetzung wurde der Pipeline-Workflow aus
Abschnitt~\ref{sec:practical:pipeline:description} um zwei Schritte erweitert:

\[
  \texttt{generate} \rightarrow \texttt{preprocess} \rightarrow \texttt{split}
  \rightarrow [\,\texttt{compute}_{1..4}\,] \rightarrow
  \texttt{aggregate} \rightarrow \texttt{postprocess}
\]

Nach dem Filtern teilt \texttt{split.py} die vorbereiteten Werte auf vier etwa gleich große
Teildateien \texttt{chunk\_1.txt} bis \texttt{chunk\_4.txt} auf. Jede dieser Teildateien wird von
einer eigenen \texttt{compute}-Instanz verarbeitet Da die vier Instanzen keine gegenseitigen
Abhängigkeiten haben, können sie parallel laufen. Der Schritt \texttt{aggregate.py} wartet auf alle
vier Teilergebnisse (Synchronization), summiert Anzahl und Summe über die Chunks und berechnet
daraus den Gesamtmittelwert. \texttt{postprocess.py} schreibt wie zuvor die abschließende
Zusammenfassung. Die zusätzlichen Skripte sowie die auf 100 erhöhte Eingabemenge sind in
Abschnitt~\ref{sec:practical:setup:scripts} beschrieben und in Anhang~\ref{app:sg:scripts}
abgebildet.

Eine Schwierigkeit beim Nachweis der Parallelität ist, dass die eigentliche Berechnung sehr kurz
ist und vier parallele Schritte daher in den Logs kaum von einer schnellen seriellen Abfolge zu
unterscheiden wären. Um die zeitliche Überlappung sichtbar zu machen, wurde in jeden
\texttt{compute}-Schritt eine künstliche Wartezeit von fünf Sekunden eingebaut außer in Nextflow. Hier ließ sich die zeitliche Überlappung dennoch zeigen. Diese Wartezeit
gehört nicht zur eigentlichen Berechnung, sondern dient als Messhilfe. Bei einer seriellen
Ausführung müsste die Compute-Phase rund zwanzig Sekunden dauern, bei echter paralleler Ausführung
nur etwa fünf bis sieben Sekunden. An diesem Unterschied lässt sich für jedes System ablesen, ob
die vier Schritte tatsächlich gleichzeitig liefen. Die folgenden Abschnitte beschreiben die
Umsetzung in den fünf Systemen; die vollständigen Logs der Läufe stehen in Anhang~\ref{app:logs},
die Abbildungen der erzeugten Reports und des DAG in Anhang~\ref{app:figures}.

\subsection{Umsetzung in Nextflow}
\label{sec:practical:scattergather:nextflow}

\paragraph{Parallel-Split.}
Der Aufbau folgt demselben DSL2-Schema wie beim Pipeline-Pattern, ergänzt um die Prozesse
\texttt{SPLIT} und \texttt{AGGREGATE} sowie einen pro Chunk aufgerufenen \texttt{COMPUTE}-Prozess.
Der Parallel-Split entsteht über zwei Channel-Operatoren. Der \texttt{SPLIT}-Prozess gibt die vier
Teildateien als eine Liste aus, und \texttt{.flatten()} zerlegt diese Liste in vier einzelne
Channel-Elemente, sodass \texttt{COMPUTE} automatisch viermal ausgeführt wird, einmal je Chunk:

\begin{lstlisting}[caption={flatten und collect im workflow-Block (Nextflow, Scatter-Gather)},label={lst:sg:nf:fc}]
chunks_ch  = SPLIT(prepared_ch).flatten()    // Parallel-Split: vier Chunks einzeln
results_ch = COMPUTE(chunks_ch).collect()    // Synchronization: Ergebnisse einsammeln
\end{lstlisting}

Für die vier parallelen Zweige war damit kein eigener Prozess nötig, ein einziger
\texttt{COMPUTE}-Prozess wird über die Operatoren mehrfach auf die Chunks angewendet. Die
vollständige Datei ist in Anhang~\ref{app:sg:nextflow:mainnf} abgebildet.

\paragraph{Synchronization.}
Die Zusammenführung übernimmt der Operator \texttt{.collect()}, der die vier Ergebnisdateien wieder
zu einer einzigen Eingabe bündelt. Erst wenn alle vier \texttt{COMPUTE}-Instanzen ein Ergebnis
geliefert haben, gibt der Channel diese gebündelte Eingabe an \texttt{AGGREGATE} weiter, sodass
\texttt{AGGREGATE} zwangsläufig auf alle vier Teilergebnisse wartet.

\paragraph{Nachweis der Parallelität.}
Dass die vier Schritte tatsächlich parallel liefen, zeigte sich zunächst an der Statusanzeige, die
für \texttt{COMPUTE} vier von vier Instanzen meldete. Über die Option \texttt{-with-timeline} erzeugt
Nextflow zusätzlich einen HTML-Zeitstrahl, in dem die vier \texttt{COMPUTE}-Aufrufe deutlich
zeitlich überlappten und \texttt{AGGREGATE} erst nach deren Abschluss begann (siehe
Anhang~\ref{app:figures}).

\paragraph{Einordnung.}
Die Umsetzung in Nextflow ist die kompakteste der fünf, da die Parallelität allein über zwei
Operatoren entsteht und eine andere Chunk-Anzahl ohne Änderung am Workflow möglich wäre. Aus
Einsteigersicht ist der \texttt{workflow}-Block sehr übersichtlich, setzt aber voraus, dass die
Bedeutung von \texttt{.flatten()} und \texttt{.collect()} verstanden wird, da sonst nicht
unmittelbar erkennbar ist, an welcher Stelle der Parallel-Split entsteht.

\subsection{Umsetzung in Pegasus}
\label{sec:practical:scattergather:pegasus}

\paragraph{Parallel-Split.}
Wie beim Pipeline-Pattern wird der Workflow über die Pegasus Python API in
\texttt{build\_workflow.py} aufgebaut. Anders als in Nextflow gibt es keinen Operator für die
Parallelität. Stattdessen werden die vier parallelen Zweige als vier separate Jobs angelegt, die
jeweils einen Chunk verarbeiten:

\begin{lstlisting}[language=Python,caption={Ein compute-Job je Chunk (Pegasus, Scatter-Gather)},label={lst:sg:pg:compute}]
compute_job_1 = Job("compute")
compute_job_1.add_inputs(chunk_1_file)
compute_job_1.add_outputs(result_1_file, stage_out=False, register_replica=False)
# compute_job_2 bis compute_job_4 analog fuer chunk_2 bis chunk_4
\end{lstlisting}

Da die vier \texttt{compute}-Jobs nur vom \texttt{split}-Job abhängen, liegen sie im DAG als
unabhängige Zweige nebeneinander und sind dadurch strukturell parallelisierbar. Die vollständige
Datei ist in Anhang~\ref{app:sg:pegasus:buildwf} abgebildet.

\paragraph{Synchronization.}
Die Zusammenführung wird über die Eingaben des \texttt{aggregate}-Jobs erreicht. Dieser nimmt alle
vier Ergebnisdateien als Eingabe entgegen:

\begin{lstlisting}[language=Python,caption={Eingaben des aggregate-Jobs (Pegasus, Scatter-Gather)},label={lst:sg:pg:agg}]
aggregate_job.add_inputs(result_1_file, result_2_file,
                         result_3_file, result_4_file)
\end{lstlisting}

Da der Job vier Eingaben hat, die von vier verschiedenen \texttt{compute}-Jobs erzeugt werden,
hängt er im DAG von allen vier ab und kann erst starten, wenn alle vier Ergebnisse vorliegen.

\paragraph{Nachweis der Parallelität.}
Zum Nachweis wurden die Start- und Endzeiten der vier \texttt{compute}-Jobs aus ihren Ausgabedateien
ausgelesen. Die Jobs starteten zwischen 14:55:59 und 14:56:01 und endeten zwischen 14:56:04 und
14:56:06. Da jeder Job künstlich fünf Sekunden gewartet hat, hätte eine serielle Ausführung rund
zwanzig Sekunden gedauert, während zwischen dem ersten Start und dem letzten Ende nur etwa sieben
Sekunden lagen. Damit ist auch bei Pegasus eine tatsächliche zeitliche Überlappung der Compute-Phase
belegt und nicht nur die strukturelle Parallelisierbarkeit, die sich bereits aus dem DAG ergibt. Die
vollständigen Logs stehen in Anhang~\ref{app:logs}, der erzeugte DAG in Anhang~\ref{app:figures}.

\paragraph{Einordnung.}
Der Modellierungsaufwand ist höher als bei Nextflow, da jeder Zweig als eigener Job definiert und
zusätzlich die Datenkonfiguration \texttt{sharedfs} sowie die Kataloge eingerichtet werden mussten.
Die Trennung von abstraktem Workflow und Ausführung ist klar gegeben, der Einstieg ist durch die
Kataloge und die Planungs- und Ausführungstrennung jedoch aufwendiger als bei Nextflow.

\subsection{Umsetzung in Merlin}
\label{sec:practical:scattergather:merlin}

\paragraph{Parallel-Split.}
Der Workflow wird wie beim Pipeline-Pattern in einer einzigen YAML-Datei
\texttt{scatter\_gather.yaml} beschrieben. Die vier parallelen Zweige wurden bei Merlin als vier eigenständige Schritte
\texttt{compute\_1} bis \texttt{compute\_4} beschrieben, die alle nur vom \texttt{split}-Schritt
abhängen und dadurch untereinander unabhängig sind:

\begin{lstlisting}[caption={Fan-out und Fan-in über depends (Merlin, Scatter-Gather)},label={lst:sg:me:depends}]
  - name: compute_1
    run:
      depends: [split]            # analog compute_2, compute_3, compute_4 
  - name: aggregate
    run:
      depends: [compute_1, compute_2, compute_3, compute_4]
\end{lstlisting}

Die vier Schritte müssen damit einzeln ausgeschrieben werden, was die Datei länger. Die vollständige Datei ist in
Anhang~\ref{app:sg:merlin:yaml} abgebildet.

\paragraph{Synchronization.}
Die Zusammenführung wird ebenfalls über \texttt{depends} erreicht. Der \texttt{aggregate}-Schritt
hängt von allen vier \texttt{compute}-Schritten ab und kann daher erst starten, wenn alle vier
abgeschlossen sind.

\paragraph{Nachweis der Parallelität.}
Der Lauf erfolgte im Worker-Modus, bei dem Merlin die Schritte als Aufgaben in Redis ablegt und von
Celery-Workern abarbeiten lässt. Die Worker-Logs zeigen, dass die vier \texttt{compute}-Schritte
praktisch gleichzeitig gestartet wurden, alle innerhalb weniger Millisekunden um 21:01:44 herum
(\texttt{compute\_4} um 21:01:44,163, \texttt{compute\_1} um 21:01:44,167, \texttt{compute\_3} um
21:01:44,170 und \texttt{compute\_2} um 21:01:44,174). Sie endeten zwischen 21:01:49,499 und
21:01:50,324, und der \texttt{aggregate}-Schritt begann erst danach um 21:01:50,365. Damit sind
beide Teile des Patterns gut sichtbar, sowohl der gleichzeitige Start der vier Schritte als auch das
Warten von \texttt{aggregate} auf alle Ergebnisse. Die vollständigen Logs stehen in
Anhang~\ref{app:logs}.

\paragraph{Einordnung.}
Beim Scatter-Gather-Pattern wird der Worker-Modus erstmals als echter Vorteil sichtbar, da die vier
Schritte tatsächlich parallel auf Worker verteilt werden. Die Verteilung über Redis und Celery
gehört zur Ausführungsumgebung, die Anzahl der parallelen Zweige ist aber fest in der Workflow-Datei
verankert. Aus Einsteigersicht bleibt die Beschreibung durch \texttt{depends} gut nachvollziehbar,
die Wiederholung der vier Schritte macht sie aber länger als bei Nextflow.

\subsection{Umsetzung in StreamFlow}
\label{sec:practical:scattergather:streamflow}

\paragraph{Parallel-Split.}
StreamFlow trennt auch hier die Workflow-Logik in \texttt{workflow.cwl} und die Ausführungsumgebung
in \texttt{streamflow.yml}. Der Parallel-Split wird über vier explizit ausformulierte Schritte
\texttt{compute\_1} bis \texttt{compute\_4} modelliert. Der \texttt{split}-Schritt stellt dazu vier
chunk Ausgaben bereit, von denen jede mit einem eigenen \texttt{compute}-Schritt verbunden ist.
Wie bei Pegasus und Merlin sind die vier Zweige damit einzeln beschrieben worden. Die
vollständigen Dateien sind in Anhang~\ref{app:sg:streamflow:cwl} abgebildet.

\paragraph{Synchronization.}
Die Zusammenführung erfolgt über den \texttt{aggregate}-Schritt, der die Ausgaben aller vier
\texttt{compute}-Schritte als Eingaben referenziert. Über diese Referenzen entsteht in der CWL die
Abhängigkeit von allen vier Zweigen, sodass \texttt{aggregate} erst nach allen vier Ergebnissen
ausgeführt wird.

\paragraph{Nachweis der Parallelität.}
Hier zeigte sich der für StreamFlow kennzeichnende Punkt. In der Standardkonfiguration startet das
Docker-Modell nur einen einzigen Container, in dem die vier \texttt{compute}-Schritte trotz
unabhängiger Modellierung nacheinander abliefen. Erst der Parameter \texttt{replicas} im Deployment
stellte mehrere Container-Instanzen bereit:

\begin{lstlisting}[caption={replicas im Deployment (StreamFlow, Scatter-Gather)},label={lst:sg:sf:replicas}]
  python-docker:
    type: docker
    config:
      image: python:3.12-slim
      replicas: 4
\end{lstlisting}

Mit \texttt{replicas: 4} verteilte StreamFlow die vier Schritte auf vier Container, und der erzeugte
HTML-Report zeigte ihre zeitliche Überlappung (siehe Anhang~\ref{app:figures}). Damit die
Überlappung deutlich wurde, wurde auch hier die künstliche Wartezeit genutzt.

\paragraph{Einordnung.}
Die Parallelität steckt bei StreamFlow nicht in der Workflow-Beschreibung, sondern in der
Ausführungsumgebung. Das passt zur klaren Trennung beider Ebenen, denn für den Wechsel von
sequenzieller zu paralleler Ausführung musste keine Zeile der CWL geändert werden, sondern nur die
Zahl der Container. Aus Einsteigersicht ist die CWL nachvollziehbar.

\subsection{Umsetzung in AiiDA}
\label{sec:practical:scattergather:aiida}

\paragraph{Parallel-Split.}
Anders als beim Pipeline-Pattern, bei dem alle Schritte Methoden einer einzigen WorkChain waren,
wurde das Scatter-Gather-Pattern auf zwei Klassen aufgeteilt. Die übergeordnete
\texttt{ScatterGatherWorkChain} steuert den Ablauf, während die eigentliche Berechnung in einer
kleineren \texttt{ComputeChunkWorkChain} steckt. Nachdem \texttt{split\_data} die vier Chunks erzeugt
hat, wird in einer Schleife für jeden Chunk mit \texttt{self.submit()} eine solche
\texttt{ComputeChunkWorkChain} an den Daemon übergeben:

\begin{lstlisting}[language=Python,caption={self.submit und ToContext (AiiDA, Scatter-Gather)},label={lst:sg:ai:submit}]
for i, chunk in enumerate(chunks, start=1):
    future = self.submit(ComputeChunkWorkChain, chunk=List(list=chunk),
                         index=Int(i), sleep=self.inputs.compute_sleep)
    futures[f"compute_{i}"] = future
return ToContext(**futures)   # warte, bis alle vier abgeschlossen sind
\end{lstlisting}

Da die vier Child-Prozesse unabhängig an den Daemon übergeben werden, können sie parallel laufen.
Die vollständige Umsetzung ist in Anhang~\ref{app:sg:aiida:workchain} abgebildet.

\paragraph{Synchronization.}
Die Zusammenführung wird über \texttt{ToContext} erreicht. Die Rückkehr aus dem Schritt mit den
gesammelten Futures bewirkt, dass die übergeordnete WorkChain so lange wartet, bis alle vier
Child-Prozesse abgeschlossen sind. Erst danach läuft sie weiter, und der Schritt \texttt{gather}
führt die vier Teilergebnisse zusammen.

\paragraph{Nachweis der Parallelität.}
Der Workflow wurde über den Daemon ausgeführt. Ein Bericht zum übergeordneten Prozess (PK 243)
zeigt, dass die vier Rechenschritte als eigene Prozesse gestartet wurden (\texttt{compute\_1} bis
\texttt{compute\_4} mit den PK 252, 255, 258 und 262), und zwar zwischen 00:46:38 und 00:46:39. Alle
vier endeten gemeinsam um 00:46:44, und die im Startskript gemessene Gesamtzeit lag bei etwa
8,7 Sekunden statt der über zwanzig Sekunden, die eine serielle Ausführung erfordert hätte. Damit
ist die zeitliche Überlappung eindeutig belegt. Die vollständigen Logs stehen in
Anhang~\ref{app:logs}.

\paragraph{Einordnung.}
Der Modellierungsaufwand war hier am größten, da neben den zusätzlichen Schritten eine zweite
WorkChain-Klasse mit eigenen Ein- und Ausgaben eingeführt werden musste. Im Gegenzug ist die
Chunk-Anzahl über einen Parameter steuerbar. Voraussetzung für den Lauf war die Ausführung über den
Daemon, der beim ersten Versuch neu gestartet werden musste, damit er den geänderten Workflow-Code
laden konnte. Wie schon beim Pipeline-Pattern speichert AiiDA auch hier die gesamte Daten- und
Prozesshistorie, sodass sich jeder Schritt im Nachhinein rekonstruieren lässt. Dieser Vorteil bei
der Nachvollziehbarkeit steht dem höchsten Aufwand bei Einrichtung und Modellierung gegenüber.

\subsection{Vergleichende Beobachtungen zum Scatter-Gather-Pattern}
\label{sec:practical:scattergather:comparison}

Anders als beim Pipeline-Pattern, das alle Systeme auf ähnliche Weise abbilden, treten beim
Scatter-Gather-Pattern größere Unterschiede zutage. Tabelle~\ref{tab:scattergather:composition}
stellt die fünf Systeme gegenüber.

\begin{table}[ht]
  \caption{Umsetzung des Scatter-Gather-Patterns je System}
  \label{tab:scattergather:composition}
  \centering
  \footnotesize
  \begin{tabular}{l p{2.5cm} p{2.5cm} p{2.7cm} p{2.6cm}}
    \hline
    \textbf{System} & \textbf{Parallel-Split} & \textbf{Gather / Sync.} & \textbf{Nachweis der Parallelität} & \textbf{Besonderheit} \\
    \hline
    Nextflow
      & \texttt{.flatten()}
      & \texttt{.collect()}
      & Timeline, \texttt{COMPUTE 4/4}
      & ein Prozess, mehrfach angewendet \\
    Pegasus
      & vier \texttt{compute}-Jobs
      & \texttt{aggregate} hängt von allen ab
      & Zeiten aus den Logs (ca. 7\,s statt 20\,s)
      & explizit, zeitlich nachgewiesen \\
    Merlin
      & vier Schritte mit \texttt{depends}
      & \texttt{aggregate} mit \texttt{depends} auf alle
      & Worker-Logs (Start fast gleichzeitig)
      & Redis und Celery, Zweige einzeln \\
    StreamFlow
      & vier CWL-Schritte
      & \texttt{aggregate} mit vier Inputs
      & HTML-Report mit \texttt{replicas: 4}
      & Parallelität in der Umgebung \\
    AiiDA
      & Schleife mit \texttt{self.submit}
      & \texttt{ToContext} und \texttt{gather}
      & \texttt{verdi}-Bericht (ca. 8,7\,s)
      & speichert Daten- und Prozesshistorie \\
    \hline
  \end{tabular}
\end{table}

Beim Ausdruck des Patterns reicht das Spektrum von einer rein operatorbasierten Modellierung
(Nextflow) über die Schleife mit \texttt{self.submit} (AiiDA) bis zur expliziten
Ausformulierung jedes Zweigs (Pegasus, Merlin, StreamFlow). Daraus folgt unmittelbar der
unterschiedliche Modellierungsaufwand: Nextflow und AiiDA kommen ohne fest eingetragene Zweige aus,
sodass sich die Chunk-Anzahl leicht ändern lässt, während sie in Pegasus, Merlin und StreamFlow im
Workflow festgeschrieben ist und bei einer anderen Aufteilung angepasst werden müsste.

Am deutlichsten unterscheiden sich die Systeme bei der Steuerung und beim Nachweis der
Parallelität. In Nextflow entsteht die Parallelität automatisch aus den Channel-Operatoren, in
Pegasus aus den Abhängigkeiten, in Merlin aus den \texttt{depends}-Beziehungen, in AiiDA aus den über den Daemon abgesetzten Child-Prozessen, und in StreamFlow erst durch die
\texttt{replicas} in der Ausführungsumgebung. Bemerkenswert ist, dass sich für alle fünf Systeme eine tatsächliche zeitliche Überlappung der Compute-Phase
nachweisen ließ, sei es über Timeline, HTML-Report, Worker-Logs, Runtime-Logs oder den
\texttt{verdi}-Bericht. 

Die Trennung von Workflow-Logik und Ausführungsumgebung ist bei StreamFlow am deutlichsten, da
dieselbe CWL je nach Deployment sequenziell oder parallel läuft und bei Pegasus durch die Trennung
von Planung und Ausführung ebenfalls klar gegeben. Bei der Nachvollziehbarkeit hebt sich erneut
AiiDA ab, das die Daten- und Prozesshistorie speichert. Aus Einsteigersicht ist Nextflow am
zugänglichsten, während die übrigen Systeme durch
Kataloge, Konfiguration, Deployment oder das AiiDA-Datenmodell höhere Hürden aufweisen.

Einschränkend ist festzuhalten, dass alle fünf Umsetzungen lokal ausgeführt wurden, also mit
lokalem Executor, lokalem Docker beziehungsweise lokalem HTCondor. Untersucht wurde damit die
Modellierbarkeit und die lokale Ausführung einer HPC-nahen Scatter-Gather-Struktur aus
Datenaufteilung, unabhängiger Verarbeitung und Zusammenführung. Die Nähe zum HPC-Umfeld ergibt sich aus der Struktur des Workflows, nicht aus
der bewusst einfach gehaltenen Rechenlogik.

\section{Systemübergreifende Erkenntnisse}
\label{sec:practical:overall}

Über beide Patterns hinweg lassen sich einige Beobachtungen zusammenfassen, die über das einzelne
System hinausgehen. Grundsätzlich konnten alle fünf Systeme beide Patterns umsetzen und
fehlerfrei ausführen, was zeigt, dass weder die sequenzielle Pipeline noch die parallele
Scatter-Gather-Struktur ein System vor unlösbare Probleme stellte. Die Unterschiede liegen weniger
darin, ob ein Pattern umsetzbar ist, sondern darin, wie aufwendig und wie nachvollziehbar die Umsetzung
ausfällt.

Nextflow kam mit beiden Patterns am unkompliziertesten zurecht. Die Beschreibung ist kompakt, die
Datenweitergabe übernimmt das System, und die Parallelität entsteht beim Scatter-Gather-Pattern
allein über zwei Operatoren, ohne dass die vier Zweige einzeln beschrieben werden müssten. Der
Preis dafür ist, dass diese Operatoren verstanden werden müssen, weil sonst nicht offensichtlich
ist, an welcher Stelle der Parallel-Split entsteht. Pegasus verlangt beim Beschreiben des Workflows
einen hohen Aufwand. Jeder Zweig muss einzeln als Job definiert, die Kataloge müssen eingerichtet und die Planung muss vor jeder Ausführung durchlaufen werden. Dafür ist der Ablauf sehr explizit, und der erzeugte DAG macht die Abhängigkeiten gut sichtbar. Beim
Scatter-Gather-Pattern ließ sich auch für Pegasus die zeitliche Überlappung über die Runtime-Logs
belegen.
 

Merlin liegt durch das Schlüsselwort \texttt{depends} gut lesbar dazwischen und das Modell aus
Redis und Celery-Workern wird beim Scatter-Gather-Pattern erstmals als echter Vorteil sichtbar, da
die vier Schritte dort tatsächlich parallel auf Worker verteilt werden. Nachteilig sind das
manuelle Kopieren der Dateien und die Notwendigkeit, die vier Zweige einzeln auszuschreiben.
StreamFlow trennt die Workflow-Beschreibung und die Ausführungsumgebung am deutlichsten, was sich
beim Scatter-Gather-Pattern besonders gezeigt hat, als für den Wechsel zu paralleler Ausführung nur
die Zahl der Container geändert werden musste. AiiDA schließlich bietet die beste Nachvollziehbarkeit, weil es als einziges System die
gesamte Daten- und Prozesshistorie speichert, erkauft sich diese Stärke aber mit dem höchsten
Aufwand bei Einrichtung und Modellierung.

Ein Punkt zeigte sich über alle Systeme hinweg und ist für den Vergleich das eigentliche Ergebnis:
Beim Scatter-Gather-Pattern ließ sich in allen fünf Fällen nicht nur eine strukturelle
Unabhängigkeit der vier Schritte, sondern eine tatsächliche zeitliche Überlappung nachweisen. Der
Weg dorthin war jedoch sehr unterschiedlich, von einem einzelnen Operator in Nextflow über vier
einzeln ausgeschriebene Zweige bei Pegasus, Merlin und StreamFlow bis zu einer Schleife mit
\texttt{self.submit} in AiiDA, und der Nachweis stützte sich je nach System auf eine Timeline,
einen HTML-Report, Worker-Logs, Runtime-Logs oder einen Prozessbericht. Dass dasselbe Muster auf so
verschiedene Arten ausgedrückt und belegt wird, verdeutlicht, wie unterschiedlich die Systeme an
dieselbe Aufgabe herangehen.

Abschließend ist die Einordnung dieser Ergebnisse wichtig. Alle Tests liefen lokal auf einem
einzelnen Rechner. Die Nähe zum HPC-Umfeld ergibt sich
allein aus der Struktur der Workflows, bei der Daten vorbereitet, aufgeteilt, unabhängig
verarbeitet und wieder zusammengeführt werden, und nicht aus einer realen Rechenlast oder einem
Scheduler, der viele Knoten verwaltet. Aus demselben Grund werden hier bewusst keine Aussagen
darüber gemacht, welches System schneller ist. Die gemessenen Zeiten sind durch den Overhead der
Systeme, durch die lokale Umgebung und nicht zuletzt durch die künstliche Wartezeit so stark
beeinflusst, dass ein Vergleich der Geschwindigkeit nicht aussagekräftig wäre. Im Mittelpunkt
stehen daher die Modellierbarkeit, die Ausführbarkeit und die Nachvollziehbarkeit der Patterns,
nicht die reine Leistung.